% Text section body for: Co od nástrojů reálně očekávat — doporučení pro pedagogy

V této části shrnujeme, co lze od generativních AI nástrojů reálně očekávat ve výuce, jak definovat „bezpečný start“ pilotu a které metriky a artefakty sledovat. Cílem je dát garantům a učitelům praktický checklist a jednoduchý workflow pro první piloty s jasnými hranicemi odpovědnosti.

Krátké očekávání schopností a omezení
\begin{itemize}
  \item Nástroje dobře pomáhají s opakovatelnými, šablonovatelnými úlohami: generování návrhů otázek, řešení příkladů, návrhů slidů nebo prvotních revizí textů. U nástrojů pro výuku matematiky například funkce „Pokrok v matematice“ generuje příklady, pomáhá s vyhodnocením postupů a připravuje návrh hodnocení, které učitel finálně schválí \cite{kb_ai_101_tipu_pdf_04e4a115_page_8_1}.
  \item Nástroje lze použít i k názorné demonstraci principů AI (např. vzdělávací světy v Minecraft Education), což je užitečné při zavádění tématu studentům \cite{kb_ai_101_tipu_pdf_04e4a115_page_15_2}.
  \item Limitace: modely mohou generovat přesvědčivě znějící, ale nepřesné nebo nepravdivé informace (tzv. „halucinace“) — toto je stručné obecné upozornění pro pedagogickou praxi (general background knowledge).
\end{itemize}

Doporučený postup „bezpečného startu“ (praktický workflow)
\begin{enumerate}
  \item Definujte rozsah a účel pilotu: konkrétní předmět, typ úloh (např. tvorba pracovních listů, formativní zpětná vazba, generování testů), časový horizont a očekávané přínosy.
  \item Proveďte DPIA a zdokumentujte datový tok: co se bude posílat do externích služeb, jaká data jsou citlivá a jak budou pseudonymizována či vyloučena; využijte vzor DPIA z příloh jako kontrolní seznam \cite{kb_malinka_chatgpt_ve_vyuce_dva_roky_pote_2024_pdf_90bc3aaf_page_12_1}.
  \item Začněte dvoufázově: fáze A bez AI (nezávislý výsledek od studentů), fáze B s AI (rozšíření, reflexe). Tím získáte srovnávací data o přínosu a rizicích nasazení \cite{kb_malinka_chatgpt_ve_vyuce_dva_roky_pote_2024_pdf_90bc3aaf_page_12_1}.
  \item Určete „human‑in‑the‑loop“ pravidla: kdo finálně schvaluje obsah, jak probíhá kontrola automaticky generovaných odpovědí a jaká je eskalace chyb.
  \item Transparentnost vůči studentům: v sylabu uveďte, kde a jak AI používáte (šablona v Příloze 9.3), a požadujte deklaraci použití tam, kde je to relevantní \cite{kb_malinka_chatgpt_ve_vyuce_dva_roky_pote_2024_pdf_90bc3aaf_page_12_1}.
\end{enumerate}

Co měřit během pilotu (praktické metriky)
\begin{itemize}
  \item Pedagogické výsledky: srovnání kvality výstupů a porozumění mezi fázemi A a B (např. hodnocení řešení, počet opravných koleček).
  \item Procesní důkazy: artefakty studentů (drafty, postupy, komentáře) a záznamy interakcí s AI (logy dotazů/odpovědí) k ověření autenticity a sledování chyb.
  \item Efektivita práce učitele: časová úspora při přípravě materiálů nebo hodnocení (např. čas na finalizaci návrhů generovaných AI), včetně kvality „předpřipravených“ návrhů, které učitel upravuje \cite{kb_ai_101_tipu_pdf_04e4a115_page_8_1}.
  \item Bezpečnostní a právní ukazatele: typy dat zasílaných do externích služeb, incidennce nežádoucích sdílení, výsledky konzultací s IT/GDPR a splnění požadavků z DPIA \cite{kb_malinka_chatgpt_ve_vyuce_dva_roky_pote_2024_pdf_90bc3aaf_page_12_1}.
\end{itemize}

Krátký checklist pro pedagoga před spuštěním pilotu
\begin{itemize}
  \item Popsaný účel pilotu a provedená DPIA (Příloha 9.4)? \cite{kb_malinka_chatgpt_ve_vyuce_dva_roky_pote_2024_pdf_90bc3aaf_page_12_1}
  \item Jsou citlivá data vyloučena nebo pseudonymizována? 
  \item Je jasně definováno, kdo finalizuje výstupy AI (human‑in‑the‑loop)? \cite{kb_ai_101_tipu_pdf_04e4a115_page_8_1}
  \item Jsou studenti informováni a je ve sylabu viditelná politika využití AI (šablona v Příloze 9.3)?
  \item Budou logovány interakce a naplánován pravidelný audit chování systému během pilotu?
\end{itemize}

Krátké příklady nasazení pro první pilot (inspirace)
\begin{itemize}
  \item Podpora přípravy: nechte AI navrhnout sadu cvičných úloh a vy jako učitel vyberete a doladíte konečný klíč — rychlý způsob, jak otestovat kvalitu návrhů bez ztráty kontroly \cite{kb_ai_101_tipu_pdf_04e4a115_page_8_1}.
  \item Didaktická demonstrace: využijte připravené světy v Minecraft Education k názornému vysvětlení, odkud model „čerpá“ a jak fungují demonstrační scénáře \cite{kb_ai_101_tipu_pdf_04e4a115_page_15_2}.
\end{itemize}

Poznámka o dalším kroku: výstupy pilotu zdokumentujte a po skončení proveďte krátké vyhodnocení pro rozhodnutí o škálování (co fungovalo, kolik práce vyžadovalo dohled, jaké úpravy v DPIA či workflow jsou nutné). Pro šablony DPIA, text do sylabu a checklisty využijte přílohy 9.3–9.5 a online repozitář ai.vut.cz jako living resource.